\section{Balance tests}
\label{sec:balance}

Retain the self-conjugate candidate $L$ in
Equation~\eqref{eq:candidate}. For this second support test we use the
full square $S^L\otimes S^L$, with both factors realized by
column-alternating tensors on the same column blocks. Put
\[
 c=2b+2,\qquad T=2M-1+r.
\]
Thus $c$ is the total size of the columns of length two, and $T$ is
the sum of the two largest column lengths. We prove that a target
occurs whenever its first four column heights have sum at most
$2M-2$. The local components below supply balanced lists of
trivial-factor sizes, whose induced support is tested by Young's
rule. We then establish an additional three-row support statement
for the bounded capacity checks.

For reference, write
\begin{align*}
 \mathcal H_j(t)&=\{\nu\vdash t:\ell(\nu)\le j\},\\
 \mathcal E_j(t)&=\{\nu\in\mathcal H_j(t):
                         \text{every row length of $\nu$ is even}\}.
\end{align*}
If $t$ is even, let $\mathcal O_4(t)$ be the members of
$\mathcal H_4(t)$ having an odd-height column.  Finally,
$\mathcal P_3(c,s)$ denotes the partitions with at most three rows
obtained from a member of $\mathcal E_3(c)$ by adding a horizontal
strip of size three and then a horizontal strip of size $s$.
All support statements below mean inclusion, without any assertion
about multiplicities.

\subsection{Separating components by output values}

Use independent ordered alphabets in the two Specht factors.
A pair of letters $(a,a')$ has output
\[
 z(a,a')=a+a'-1.
\]
For a collection $B$ of column blocks on a position set $D_B$, let
$v_B$ be their alternating tensor and put
$Y_B=\C[S_{D_B}]v_B$. By Equation~\eqref{eq:block-specht}, this is
the Specht module whose column lengths are the sizes of the blocks
in $B$. Realize both copies of $Y_B$ on $D_B$. Every pair word
occurring in their tensor square has the same number of positions
and the same output sum:
\begin{equation}
 d_B=\sum_{A\in B}|A|,
 \qquad
 \sigma_B=\sum_{A\in B}|A|^2.
 \label{bal:fixed-sums}
\end{equation}
Indeed, a block of length $k$ uses $1,\ldots,k$ once in each
factor, so its output sum is $2(1+\cdots+k)-k=k^2$. Permuting
positions within $D_B$, independently in the two layers, preserves
these totals.

\begin{lemma}[Separation by sizes and sums]
\label{lem:balance-separation}
Partition the column blocks of $L$ into collections
$B_1,\ldots,B_v$ on disjoint position sets $D_1,\ldots,D_v$,
using the same blocks in both factors. Suppose $I_1,\ldots,I_v$ are
nonempty sets of integers with every element of $I_i$ smaller than
every element of $I_{i+1}$. Let $W_i$ be the image of
$Y_{B_i}\otimes Y_{B_i}$ under the coordinate projection retaining
words all of whose outputs lie in $I_i$. Then the support of
$S^L\otimes S^L$ contains the support of
\[
 \Ind_{S_{d_{B_1}}\times\cdots\times S_{d_{B_v}}}^{S_n}
       (W_1\boxtimes\cdots\boxtimes W_v).
\]
\end{lemma}

\begin{proof}
In the full pair-word space, retain precisely the words whose
outputs lie in $\bigcup_i I_i$ and for which the outputs in $I_i$
have cardinality $d_{B_i}$ and sum $\sigma_{B_i}$, for every $i$.
This coordinate projection commutes with all position permutations.

Consider a word from the tensor product with the prescribed
component positions. The positions in $D_1$ have
cardinality $d_{B_1}$ and sum $\sigma_{B_1}$ by
Equation~\eqref{bal:fixed-sums}.  In a retained word, the positions
whose outputs lie in $I_1$ have the same cardinality and sum.
If these two sets of positions differed, the positions in the first
set but not the second would all have larger outputs than the
positions in the second but not the first.  The two set differences
have equal positive cardinality, contradicting equality of their
sums.  The two sets therefore agree.  Remove them and repeat the
argument for $I_2$, and then successively for every range.

It follows that the global projection on these prescribed positions
is the tensor product of the local projections. More explicitly,
write $v_i=v_{B_i}$, $Y_i=Y_{B_i}$, and
$H=\prod_i S_{D_i}$. Inside either Specht factor,
\[
 \C[H](v_1\otimes\cdots\otimes v_v)
       =Y_1\boxtimes\cdots\boxtimes Y_v\subseteq S^L.
\]
Taking its tensor square gives the actual subspace
$\boxtimes_i(Y_i\otimes Y_i)$ in $S^L\otimes S^L$: the two
layers can be generated independently. The symmetric groups act
diagonally on each local square and on the global square. The
restriction of the global projection is $\bigotimes_i P_i$, where $P_i$ is the
local projection.  Its image is therefore the full outer product
$W_1\boxtimes\cdots\boxtimes W_v$.
Different ordered placements of these components have disjoint
coordinate supports, distinguished by which positions have outputs
in each $I_i$.  Their translates therefore give the full induced
module in the statement.  Since an equivariant image cannot acquire
a new irreducible type, the result follows.
\end{proof}

The local support guarantees proved in the following subsections are
summarized in Table~\ref{bal:local-table}. The entries concerning the largest
two lengths use $q=M+b-1$.  When $s=1$, their lengths are $q+1,q$;
when $s=2$, their lengths are $q+2,q$.

\begin{table}[htbp]
\centering
\small
\renewcommand{\arraystretch}{1.25}
\begin{tabular}{@{}p{0.25\linewidth}p{0.15\linewidth}p{0.19\linewidth}p{0.29\linewidth}@{}}
\hline
Column blocks & Total size & Allowed outputs
    & Guaranteed support\\
\hline
$a$ copies of $k$ & $ak$ & $\{k\}$ & $(ak)$\\
$u,u-1$ & $2u-1$ & $[u-1,u]$ & $\mathcal H_4(2u-1)$\\
$q+2,q$ & $2q+2$ & $[q,2q+3]$ & $\mathcal O_4(2q+2)$\\
$b+1$ copies of $2$ & $c$ & $\{2\}$ & $\mathcal E_2(c)$\\
$3$, $b+1$ copies of $2$, and $s$ copies of $1$
    & $c+3+s$ & $[1,3]$ & $\mathcal P_3(c,s)$\\
\hline
\end{tabular}
\caption{Local projected supports.  Neighboring intervals in a
packing are chosen disjoint before applying
Lemma~\ref{lem:balance-separation}.}
\label{bal:local-table}
\end{table}

After applying Lemma~\ref{lem:balance-separation}, transitivity
of induction permits us to combine selected factors in any order.
In particular, the big pair can be combined with the singleton
component even though their output ranges are not adjacent.

\subsection{Coordinate functionals and path components}

We first explain exactly how a surviving word coordinate supplies
a path vector after the two sign twists.

\begin{lemma}[The sign-twisted coordinate line]
\label{bal:coordinate-line}
Let $X=S^\mu$ be realized by column-alternating vectors in the word
space of content $\mu=(\mu_1,\mu_2,\ldots)$.  For a word $a$ of
this content, let $R_a$ permute separately the positions carrying
each equal letter.  The restriction $\epsilon_a|_X$ of its
coordinate functional is nonzero.  Under
$X^*\otimes\sgn\simeq S^{\mu^t}$, its line is the
line of the column-alternating vector on the equal-letter blocks
of $a$, whose lengths are $\mu_1,\mu_2,\ldots$.
\end{lemma}

\begin{proof}
For the row-constant word of a tableau, its coefficient in the
tableau's column-alternating vector is one: only the identity in
each column stabilizes that word.  Hence its coordinate restriction
is nonzero.  Every other word of the same content is a position
translate, which proves nonvanishing in general.

The restricted coordinate is $R_a$-invariant. Young's rule and
Frobenius reciprocity give
\[
 \dim (X^*)^{R_a}=1,
\]
using the diagonal multiplicity in Lemma~\ref{lem:young-support}.
After twisting by sign, this becomes the unique line on which
$R_a$ acts by its sign character.  The nonzero alternating vector
on the equal-letter blocks is such a vector, and its column
lengths are the parts of $\mu$; it therefore belongs to
$S^{\mu^t}$.  This identifies the two lines.
The identification is equivariant, so translating a word changes
the identification only by the corresponding sign and the
position translate.
\end{proof}

In a tensor square, these two sign twists cancel. To make the
passage through a projection explicit, let
$P:X\otimes X\longrightarrow W$ be a local coordinate projection
onto its image. Its dual embeds $W^*$ in $(X\otimes X)^*$.
For a surviving pair word, its coordinate functional $\epsilon$
satisfies $\epsilon(Pz)=\epsilon(z)$, so its restriction belongs
to this embedded copy of $W^*$. Lemma~\ref{bal:coordinate-line}
identifies it with a nonzero multiple of the product of the two
alternating vectors on its equal-letter blocks. Its entire
position-permutation span remains in $W^*$. Finding a path among
these products therefore proves support in $W$ itself, since
complex Specht modules are self-dual. In the next two lemmas, zero
column lengths are omitted.

\begin{lemma}[Neighboring lengths]
\label{bal:neighboring}
For every integer $u\ge1$, the local component with column lengths
$u,u-1$, projected to
outputs in $[u-1,u]$, contains every type in
$\mathcal H_4(2u-1)$.
\end{lemma}

\begin{proof}
Its content is $(2^{u-1},1)$.  The equal-letter blocks thus consist
of $u-1$ pairs and one singleton in each layer.  Arrange the two
pairings as one odd path.  From its first-layer singleton, label
that singleton $u$, and label the successive first-layer pairs
$u-1,u-2,\ldots,1$.  Label the second-layer pairs, in path order,
$1,2,\ldots,u-1$, and label its terminal singleton $u$.
At the path positions the outputs are alternately $u$ and $u-1$,
with output $u$ at both ends.

The pair word survives the stated projection.  By
Lemma~\ref{bal:coordinate-line}, its coordinate functional gives
the odd-path tensor in the sign-twisted dual square.  The twists
cancel, and Lemma~\ref{band:path-support} gives every partition
with at most four rows.  Consequently all these types occur in
the projected component.
\end{proof}

\begin{lemma}[The two-path component]
\label{bal:two-paths}
For every integer $q\ge0$, the component with column lengths
$q+2,q$, projected to outputs
in $[q,2q+3]$, contains $\mathcal O_4(2q+2)$.  If $q\ge3$,
inducing this support together with a trivial representation of
size two contains every member of $\mathcal H_4(2q+4)$.
\end{lemma}

\begin{proof}
The content is $(2^q,1,1)$.  In each layer there are pair labels
$1,\ldots,q$ and singleton labels $q+1,q+2$.
For any $0\le p\le q$, make odd paths of sizes $2p+1$ and
$2(q-p)+1$.  On the first path use first-layer pair labels
$1,\ldots,p$ and second-layer pair labels $q-p+1,\ldots,q$;
on the second path use the complementary intervals.  In path
order, the first-layer pairs decrease and the second-layer pairs
increase.  Give the first path the initial first-layer singleton
$q+1$ and terminal second-layer singleton $q+2$, and the second
path the remaining singletons.

Internal outputs are $q$ or $q+1$.  Endpoint outputs are at least
$q$, and all outputs are at most $2q+3$.  If one path is a
singleton, its two singleton labels satisfy the same bounds.
The resulting pair word therefore survives the projection, and
Lemma~\ref{bal:coordinate-line} identifies its coordinate with
the product of the two path tensors.

Let $\nu\in\mathcal O_4(2q+2)$.  Since its size is even, it has
an even positive number of odd-height columns.  Divide its columns
into two groups, each of odd total size.  These totals are
$2p+1$ and $2(q-p)+1$ for an allowed $p$.  Arrange each group's
columns consecutively on the corresponding path.  By
Lemma~\ref{band:path-support}, each contraction is a nonzero
polynomial in a common ordered basis of the pair-letter dual
space. Their product is also nonzero,
so a common generic ordered basis gives a nonzero contraction of the
two-path tensor with the dual tableau of $\nu$.  This proves the
first assertion.

For the second assertion, let $\eta\in\mathcal H_4(2q+4)$.
It has width greater than two when $q\ge3$.  We claim that some
removable horizontal two-strip leaves an odd-height column.
If $\eta$ has no odd-height columns, removing any horizontal
two-strip creates two.  If it has at least four, such a strip
cannot remove all of them.  If it has exactly two, its width
ensures that an even-height column also exists.  Remove one cell
from an odd-height column and one from an even-height column,
choosing the rightmost column of each selected height.  Their
heights are different, and these removals preserve the weakly
decreasing order of column heights.  This is a horizontal
two-strip and leaves two odd-height columns.  Pieri now proves
the desired containment.
\end{proof}

For the candidate diagram, the big component has $T$ positions.
When $s=1$, Lemma~\ref{bal:neighboring} supplies
$\mathcal H_4(T)$.  When $s=2$, its two singleton columns can
supply a trivial representation of size two at output one;
Lemma~\ref{lem:balance-separation} and
Lemma~\ref{bal:two-paths} then supply $\mathcal H_4(T+2)$.
The big component's outputs are at least $M-1$ in both cases.

\subsection{Symmetric pairs and short strips}

The path calculations handle the long columns. For the short-column
components we use a direct fact about repeated symmetric tensors.
Its polynomial-module conclusion is Littlewood's classical even-row
decomposition of symmetric powers of a symmetric square
\cite[Chapter~I, Section~5, Example~5(a); Appendix~A, Section~7,
Example~(1)]{Macdonald}.
The explicit contraction below also detects the prescribed tensor,
which is what our coordinate projection requires.

\begin{lemma}[Repeated symmetric pairs]
\label{bal:even-rows}
Let $m\ge0$ be an integer, let $E$ have dimension $j$, and let
$H\in\Sym^2 E$ be
nondegenerate.  The position-permutation span of
$H^{\otimes m}\in E^{\otimes 2m}$ contains every member of
$\mathcal E_j(2m)$.  Moreover, the polynomial $\GL(E)$-module
$\Sym^m(\Sym^2 E)$ contains the Schur module of each such
partition.
\end{lemma}

\begin{proof}
Choose an orthonormal basis for $H$, so
$H=\sum_{i=1}^j e_i\otimes e_i$.  A diagram with even row lengths
has its columns paired into equal-height pairs.  Place the two
positions of each copy of $H$ in corresponding cells of two
paired columns.  Contract against the usual column-alternating
dual tensor.  For a pair of columns of height $h$, the two
permutations must agree, their signs multiply to one, and their
contribution is $h!$.  The entire contraction is the nonzero
product of these factorials.  Lemma~\ref{lem:dual-test} proves
the first statement.

For the second, use the natural embedding of
$\Sym^m(\Sym^2 E)$ into $E^{\otimes 2m}$ with the $m$ pairs
of positions fixed, and choose the dual tableau placements as
above.  The same contraction proves that this $\GL(E)$-stable
subspace has a nonzero projection to the indicated Schur-Weyl
isotypic space.  Complete reducibility then gives the required
Schur module.
\end{proof}

An isolated length-$k$ block, projected to constant output $k$,
has the pair labels $(a,k+1-a)$ for $1\le a\le k$.  Reversing
one block's labels shows that all terms in the projected tensor
have the same sign.  It is a nonzero symmetric tensor and hence
supplies the trivial type of size $k$.  For several equal-length
blocks, the projected product still has coefficients of one
common sign.  Its full position average is nonzero, supplying
the trivial type of their total size.  This proves the first
row of Table~\ref{bal:local-table}.

On an aligned length-two block, projection to constant output two
gives, up to a common nonzero scalar,
\[
 e_{12}\otimes e_{21}+e_{21}\otimes e_{12}.
\]
This is a nondegenerate symmetric form on the two-letter space
$\Span(e_{12},e_{21})$.  Thus $b+1$ such blocks supply
$\mathcal E_2(c)$ by Lemma~\ref{bal:even-rows}.

\begin{lemma}[An odd horizontal strip]
\label{bal:odd-strip}
Let $c$ be positive and even, and let $x\in\{1,3\}$ with
$x\le c$.  Inducing the support $\mathcal E_2(c)$ with the
trivial representation of size $x$ supplies
$\mathcal H_2(c+x)$.
\end{lemma}

\begin{proof}
Write a target as $(\alpha,\beta)$, allowing $\beta=0$.
A precursor of size $c$ with even row lengths has the form
$(c-y,y)$ with $y$ even.  The horizontal-strip interlacing
conditions are precisely
\begin{equation}
 \max(0,\beta-x)\le y\le\min(\beta,c-\beta).
 \label{bal:parity-interval}
\end{equation}
The upper bound is at most $c/2$, so this also guarantees that
the precursor is a partition.  If the interval contains zero,
choose $y=0$.  Otherwise its lower endpoint is $\beta-x>0$.
When $\beta\le c/2$, its length is $x\ge1$.  When
$\beta>c/2$, its length is
$c+x-2\beta=\alpha-\beta$, a positive integer because
$c+x$ is odd.  The interval therefore contains an even integer.
Pieri gives the assertion.
\end{proof}

In our applications, the size-one strip is a singleton column
at output one, and the size-three strip is the length-three
block at output three.  These ranges are disjoint from output
two, so Lemma~\ref{lem:balance-separation} justifies the required
inductions.

\subsection{The first-four-column bound}

We now combine the local path and symmetric-pair supports. The aim is
to replace each component by a list of trivial-factor sizes and then
verify Young's dominance inequalities for the entire target.

For positive integers $t,j$, let $\operatorname{pack}(t,j)$
be the list of $j$ integers whose entries differ by at most one
and sum to $t$; zero entries, if any, are omitted.  Young's rule
implies that the support obtained by inducing trivial factors of
these sizes is exactly $\mathcal H_j(t)$.  Indeed, every partition
of $t$ with at most $j$ rows dominates the balanced $j$-part
partition, whereas a partition with more than $j$ rows fails
its $j$th dominance inequality.  Thus a local component supplying
$\mathcal H_j(t)$ may contribute $\operatorname{pack}(t,j)$
to a list of trivial-factor sizes for a support calculation.

\begin{proposition}[The balance bound]
\label{prop:balance-four}
Assume $M\ge9$, and let $h_1\ge h_2\ge\cdots$ be the column
heights of $\nu\vdash n$.  If
\begin{equation}
 \sum_{i=1}^4 h_i\le2M-2,
 \label{bal:four-small}
\end{equation}
then $S^\nu$ occurs in $S^L\otimes S^L$.
The same conclusion holds if Equation~\eqref{bal:four-small}
holds for the column heights of $\nu^t$.
\end{proposition}

\begin{proof}
Use the big pair, together with the two singleton columns when
$s=2$, to contribute four balanced parts of total $T+2(s-1)$.
When $s=1$, retain its singleton column as a separate part one.
If $c\le M$, use the length-two columns to contribute the single
part $c$.  Pair the middle lengths $3,\ldots,M-2$ consecutively,
leaving their smallest member alone if their number is odd.
Each paired pair of neighboring lengths contributes four balanced
parts by Lemma~\ref{bal:neighboring}, and an unpaired length
contributes itself.

If $c>M$, use the length-three block with the length-two blocks.
Lemma~\ref{bal:odd-strip} contributes two balanced parts of total
$c+3$, and now pair the remaining middle lengths $4,\ldots,M-2$
in the same fashion.  All output ranges used are disjoint:
the singletons have output one, the length-two and length-three
constructions use outputs two and three, a middle neighboring
pair uses the interval between its two lengths, and the big
pair uses outputs at least $M-1$.  Apply
Lemma~\ref{lem:balance-separation} first to these separate
components.  Then reassociate induction to combine the big
pair with the output-one singleton component when $s=2$, and
to combine the output-two and output-three components when
$c>M$.  The local support statements above now validate the
specified packing; no disconnected output range is treated
as a single component in the separation lemma.
Call the resulting list of parts $A$; its sum is $n$.

We record two uniform bounds:
\begin{equation}
 |A|\ge2M-5,
 \qquad \max A\le M+1.
 \label{bal:packing-bounds}
\end{equation}
A run of $p$ middle lengths contributes
$2p-(p\bmod2)\ge2p-1$ parts.  If $c\le M$, there are
$p=M-4$ middle lengths and at least five other parts, giving
at least $2M-4$ parts in total.  If $c>M$, then $M$ must be
even: for odd $M$, the bound $r\le M/2$ gives
$c\le r+2\le M/2+2<M$.  In the even case, $p=M-5$ is odd,
so the middle contribution is $2M-11$, and the big and short
components contribute at least six more parts.  This proves
the first bound.

Untouched lengths and balanced middle parts are at most $M$.
The largest big-component part is at most
\[
 \left\lceil\frac{T+2}{4}\right\rceil
 \le \left\lceil\frac{7M/2+3}{4}\right\rceil
 \le M+1,
\]
using $r\le3M/2+2$.  In the exceptional case $c>M$, the
largest short-component part is
\[
 \left\lceil\frac{c+3}{2}\right\rceil
   =\left\lfloor\frac r2\right\rfloor+3
   \le\left\lfloor\frac{3M}{4}\right\rfloor+4
   \le M+1.
\]
The last inequality holds for even $M\ge10$: it is equality
at $M=10$, and follows from $3M/4+4\le M+1$ for $M\ge12$.
This proves the second bound in
Equation~\eqref{bal:packing-bounds}.  In particular,
\begin{equation}
 \frac{n}{\max A}\ge\frac{N_M}{M+1}
      =\frac M2>\frac{2M-2}{4}.
 \label{bal:average-bound}
\end{equation}

We verify the inequalities of Lemma~\ref{lem:young-support}.
The target has at least four columns, since otherwise
Equation~\eqref{bal:four-small} would give
$n\le2M-2<N_M$.  Consequently $h_1\le2M-5$, and the first
inequality follows from $P_A(1)=|A|$.  The four big-component
parts are each at least two, so
\[
 P_A(2)\ge |A|+4\ge2M-1.
\]
This proves the inequalities for prefixes two and three as well.

For $k\ge4$, decreasing column heights and
Equation~\eqref{bal:four-small} imply
\[
 \sum_{i=1}^k h_i\le
       \min\left(n,\frac{k(2M-2)}4\right).
\]
Writing $B=\max A$, the termwise bound
$\min(k,a)\ge ka/B$ for $k\le B$ gives
\[
 P_A(k)\ge \min\left(n,\frac{kn}{B}\right).
\]
Equation~\eqref{bal:average-bound} proves every remaining
dominance inequality.  The constructed support therefore contains
$S^\nu$.  Finally, since $L$ is self-conjugate, its square
is invariant under sign twist, so the same conclusion follows
from the condition on $\nu^t$.
\end{proof}

\subsection{A three-row component}

The following construction is independent of the big pair and
will be used in the cases $M=10,12,14$.

\begin{proposition}[The short three-row component]
\label{prop:three-row}
Group the length-three block, the $b+1$ length-two blocks, and
the $s$ singleton blocks in both layers, and project to outputs
in $[1,3]$.  The resulting module contains every partition in
$\mathcal P_3(c,s)$.
This component can be induced independently with components on
the remaining lengths $4,\ldots,M-2$ and with the big pair.
\end{proposition}

\begin{proof}
Align corresponding blocks in the two layers and use pair letters
$e_{aa'}$.  Set
\[
 p=e_{11},\quad q=e_{22},\quad u=e_{12},\quad v=e_{21},
 \qquad a=e_{13},\quad d=e_{31}.
\]
These six pair letters form a basis of the space of outputs at
most three.  Each aligned length-two block contributes
\begin{equation}
 H=p\otimes q+q\otimes p-u\otimes v-v\otimes u.
 \label{bal:four-form}
\end{equation}
Each singleton contributes $p$.  The length-three block has
fixed output sum nine; since all its outputs are at most three,
they must all equal three.  Its contribution is, up to a nonzero
scalar, the full symmetrization of $a,d,q$.  Thus the projected
module contains the position orbit of
\begin{equation}
 G=H^{\otimes(b+1)}\otimes p^{\otimes s}
       \otimes\sum_{\sigma\in S_3}
                z_{\sigma(1)}\otimes z_{\sigma(2)}
                              \otimes z_{\sigma(3)},
 \qquad (z_1,z_2,z_3)=(a,d,q).
 \label{bal:three-generator}
\end{equation}

We want to vary the symmetric-pair, singleton, and triple factors
independently so that Pieri applies to their tensor product.
Let $P=\Span(p,q,u,v,a,d)$. Common linear maps from $P$ provide
the needed vectors without adding any symmetric-group types.
More precisely, let $E=\C^3$, let $D=c+s+3$, and set
\[
 N=\C[S_D]G,\qquad
 W=\sum_{\phi:P\to E}\phi^{\otimes D}(N)
       \subseteq E^{\otimes D}.
\]
Each $\phi^{\otimes D}$ intertwines the position action, so every
$S_D$-type in $W$ already occurs in $N$. Moreover, $W$ is
$\GL(E)$-stable, since postcomposing $\phi$ with an element of
$\GL(E)$ gives another map in the sum. The choice $\dim E=3$
restricts the Schur-Weyl types to partitions with at most three rows.

Our immediate aim is to put
\begin{equation}
 Q^{\otimes(b+1)}\otimes x^{\otimes s}\otimes y^{\otimes3}
 \label{bal:independent-tensors}
\end{equation}
in $W$ for generic independently chosen symmetric tensors
$Q\in\Sym^2 E$ and vectors $x,y\in E$. Here $Q$ will be the
image of $H$, while $x$ and $y$ control the singleton and triple
factors. The same source letter $q$ occurs in both $H$ and the
triple, so these choices require a construction.

Identify $Q$ with its symmetric matrix and assume it is invertible.
We seek a map with $p\mapsto x$ and $q,a,d\mapsto ty$ for some
$t\ne0$. The requirement $H\mapsto Q$ then asks that
$t(xy^{\mathsf T}+yx^{\mathsf T})-Q=UV^{\mathsf T}+VU^{\mathsf T}$
for the still unchosen images $U,V$ of $u,v$. Since the right side
has rank at most two, we choose $t$ to make the left side singular.
Put
\[
 \alpha=x^{\mathsf T}Q^{-1}y,
 \qquad \beta=x^{\mathsf T}Q^{-1}x,
 \qquad \gamma=y^{\mathsf T}Q^{-1}y.
\]
The determinant lemma gives
\begin{equation}
 \det\bigl(Q-t(xy^{\mathsf T}+yx^{\mathsf T})\bigr)
   =\det(Q)\bigl(1-2\alpha t+
                         (\alpha^2-\beta\gamma)t^2\bigr).
 \label{bal:density-determinant}
\end{equation}
The quadratic coefficient is generically nonzero: for $Q=I$,
$x=e_1$, and $y=e_2$, it equals $-1$. The constant coefficient
is one. Over $\C$, there is therefore a
nonzero root $t$.  The symmetric matrix
\[
 R=t(xy^{\mathsf T}+yx^{\mathsf T})-Q
\]
has rank at most two.  Every such symmetric matrix over $\C$
can be written $UV^{\mathsf T}+VU^{\mathsf T}$: diagonalize it
by congruence, write it as $AA^{\mathsf T}+BB^{\mathsf T}$
with at most two summands, and take
$U=(A+\mathrm i B)/\sqrt2$ and
$V=(A-\mathrm i B)/\sqrt2$.

Consequently the common linear map $\phi:P\to E$ given by
\begin{equation}
 p\longmapsto x,\quad q\longmapsto ty,\quad
 u\longmapsto U,\quad v\longmapsto V,\quad
 a\longmapsto ty,\quad d\longmapsto ty
 \label{bal:density-map}
\end{equation}
sends $H$ to $Q$ and the symmetric triple to
$6t^3y^{\otimes3}$. Its image of $G$ is therefore $6t^3$ times
Equation~\eqref{bal:independent-tensors}. Since $t\ne0$ and $W$
is a linear space, the desired tensor belongs to $W$.

Varying the generic choices of $Q,x,y$ shows that $W$ contains
the naturally embedded space
\begin{equation}
 \Sym^{b+1}(\Sym^2 E)\otimes\Sym^s E\otimes\Sym^3 E.
 \label{bal:pieri-space}
\end{equation}
Indeed, a linear functional annihilating $W$ gives a
multihomogeneous polynomial in $(Q,x,y)$ when evaluated on
Equation~\eqref{bal:independent-tensors}.  It vanishes on a dense
set and hence identically.  Pure powers span each symmetric
power, proving the containment in
Equation~\eqref{bal:pieri-space}.

By Lemma~\ref{bal:even-rows}, the first factor contains the
Schur module for every member of $\mathcal E_3(c)$.  Two
applications of Pieri give every Schur module indexed by
$\mathcal P_3(c,s)$.  Schur-Weyl duality transfers these
$\GL(E)$-types to the corresponding $S_D$-types in $W$, hence
in $N$ and in the local projected module.  This proves the
asserted support.

Finally, the remaining middle components can use outputs at least
four and at most $M-2$, while the big pair uses outputs at least
$M-1$.  They are disjoint from $[1,3]$, so their independence
follows from Lemma~\ref{lem:balance-separation}.
\end{proof}

\subsection{Using the local supports in capacity tests}

The capacity tests use these local supports in two ways. First,
they retain partitions with at most three rows whose entire dominance upper
set lies in $\mathcal P_3(c,s)$.  The precise condition in
Equation~\eqref{eq:capacity-low-mask} guarantees the full
induced-trivial support associated with each retained partition's
row lengths.  Combining these lower supports with balanced middle
parts gives the dominance bounds used below.

Second, they allocate at most four cells from each target column
to the big pair, with total allocation $T$ and at least one odd
allocation when $s=2$.  If the residual diagram is supplied by
the other components, Lemma~\ref{lem:lr-addition} combines the
columnwise allocations to give the target.  The exact allocation
test and its support justification are given in
Lemma~\ref{lem:capacity-hit}.
